\subsection{Measure associated with a twisted differential form}

Recall that, in this section, we assume K = R and that all manifolds considered are locally finite-dimensional.

10.4.1. Let X be a manifold of class \(C^{r}\). One can apply the definitions and results of 10.3 by taking as vector bundle M the tangent bundle T(X). One then has \(\mathrm{Or}_{\mathbf{M}} = \tilde{\mathbf{X}}\) (10.2.4); we denote it by \(\tilde{\mathbf{R}}_{\mathbf{X}}\) and call simply the bundle of twisted scalars the bundle \(\tilde{\mathbf{R}}_{\mathrm{T}(\mathbf{X})}\) (10.3.1)\footnote{Note that this notation conflicts with the notation \(\widetilde{\mathbf{R}}_{\mathbf{M}}\) when X is regarded as endowed with its vector-bundle structure of rank 0 over itself.}; a T(X)-twisted differential form with values in a vector bundle F is called simply a twisted (or odd) differential form with values in F. When X is equipped with an orientation, the map \(\omega \mapsto \xi \otimes \omega\) allows one to identify the usual differential forms (sometimes called "even") with the twisted forms.

10.4.2. Let \(f\colon\mathbf{X}\to\mathbf{Y}\) be a morphism of manifolds of class \(\mathbf{C}^{r}\) and let \(\tilde{f}\colon\tilde{\mathbf{X}}\to\tilde{\mathbf{Y}}\) be an orientation of \(f\) (10.2.5). There exists one and only one isomorphism \(j\) from \(f^{*}(\widetilde{\mathbf{R}}_{\mathbf{Y}})\) onto \(\widetilde{\mathbf{R}}_{\mathbf{X}}\) such that \(\tilde{x}=j(\pi(\tilde{x}),\tilde{f}(\tilde{x}))\) for all \(\tilde{x}\in\tilde{\mathbf{X}}\). We identify these two bundles by means of \(j\). If \(\omega\) is a twisted form of degree \(p\) on Y with values in a vector bundle F, the inverse image \(f^{*}(\omega)\) is identified with a twisted differential form of degree \(p\) on X with values in \(f^{*}(\mathrm{F})\). If one denotes by \(\tilde{\omega}\) (resp. \(\widetilde{f^{*}(\omega)}\)) the differential form on \(\tilde{\mathbf{Y}}\) (resp. \(\tilde{\mathbf{X}}\)) corresponding to \(\omega\) (resp. \(f^{*}(\omega)\)) as in 10.3.3, one has \(\tilde{f}^{*}(\tilde{\omega})=\widetilde{f^{*}(\omega)}\).

When F is the trivial bundle defined by a Banach space, the operation \(f^*\) commutes with exterior differentiation (10.3.4).

10.4.3. Let X be a pure manifold of dimension n, and let \(\omega\) be a twisted differential form of degree n on X, with values in a Banach space E. Let \(c = (\mathrm{U}, \varphi, \mathbf{R}^{n})\) be a chart of X and let \(\xi^{n}\) be the canonical orientation of \(R^{n}\); denote by \(u^{1}, \ldots, u^{n}\) the coordinate functions on \(R^{n}\). There exists one and only one function \(f_{c}\) on \(\varphi(\mathrm{U})\), with values in E, such that

\[
\omega | \mathrm{U} = \varphi^ {*} (\xi^ {n} \otimes f _ {c}. d u ^ {1} \wedge \dots \wedge d u ^ {n}).
\]

One says that \(\omega\) is locally integrable if, for every chart \(c = (\mathrm{U},\varphi ,\mathbb{R}^n)\) of X, the corresponding function \(f_{c}\) is locally integrable with respect to \(\lambda_{\varphi (\mathrm{U})}^{\otimes n}\) (where \(\lambda\) denotes Lebesgue measure on R). Suppose that this is the case and that X is separated. There exists on X a vector measure \(\alpha (\omega)\) with values in E and one and only one possessing the following property: for every chart \(c = (\mathrm{U},\varphi ,\mathbb{R}^n)\) of X, the image under \(\varphi\) of the restriction to U of the measure \(\alpha (\omega)\) is the measure \(f_{c}\lambda_{\varphi(\mathrm{U})}^{\otimes n}\) (INT, VI, § 2, no. 4). One says that \(\alpha (\omega)\) is the measure defined by \(\omega\) and one denotes it most often simply by \(\omega\). The support of \(\alpha (\omega)\) is contained in the support \(\operatorname{Supp}\omega\) of \(\omega\) (calling the support of a differential form \(\omega\) the closure of the set of \(x\in X\) such that \(\omega (x)\neq 0\)).

If \(f\) is a real function on X, locally integrable with respect to the measure \(\alpha(\omega)\), then the differential form \(f\omega\) is locally integrable and one has \(\alpha(f\omega) = f.\alpha(\omega)\).

If \(g\) is a real function on X essentially integrable for \(\alpha(\omega)\), its integral is generally denoted \(\int_{\mathbf{X}} g\omega\), or \(\int g\omega\), or \(\int g(x)\omega(x)\); one avoids denoting it \(\int g d\omega\), because of the risk of confusion with the exterior differential \(d\omega\) of \(\omega\), which is defined (and moreover equal to 0) if \(\omega\) is of class \(C^{1}\). If A is a subset of X whose characteristic function \(\varphi_{\mathrm{A}}\) is essentially integrable for \(\alpha(\omega)\), the integral of \(\varphi_{\mathrm{A}}\) is denoted \(\int_{A} \omega\). If the constant 1 is essentially integrable for \(\alpha(\omega)\), one says that \(\omega\) is integrable.

Now let \(p\) be an integer, with \(0 \leqslant p \leqslant n\), and let \(\omega\) be a twisted differential form of degree \(p\) on X, with values in a Banach space E. Let moreover Y be a pure submanifold of dimension \(p\) of X; suppose that the canonical injection \(i: Y \to X\) is equipped with an orientation. The inverse image \(i^{*}(\omega)\) (10.4.2) is then sometimes denoted \(\omega|Y\). If \(\omega|Y\) is integrable, we set

\[
\int_ {\mathbf {Y}} \omega = \int_ {\mathbf {Y}} \omega | \mathbf {Y}.
\]

10.4.4. Let X be a pure manifold of dimension n, equipped with an orientation \(\xi\). The identification \(\omega \mapsto \xi \otimes \omega\) between usual differential forms and twisted differential forms makes it possible to apply the preceding results to forms of degree n on X; thus, to every locally integrable form \(\omega\) of degree n, there is associated a measure on X, still denoted \(\omega\). If E = R, such a form \(\omega\) is locally of integrable modulus (10.1.5) and the measure \(\text{mod}(\omega)_{\mu}\) corresponding to it (10.1.6) is none other than the absolute value \(|\omega|\) of the measure \(\omega\) (INT, III, § 1, no. 6).

10.4.5. Example. --- Let X be a pure real manifold of dimension 1, oriented, separated, and let \(z: \mathrm{X} \to \mathbf{C}\) be a differentiable map from X to C. The complex differential form \(dz\) defines on X a complex measure also denoted \(dz\) . This applies in particular when X is a pure differentiable submanifold of dimension 1 of C, \(z\) being the injection of X into C.

More specifically, let X be a circle with center 0 and radius \(\rho > 0\); orient X as indicated in 10.2.8, b), given the usual identification of C with \(\mathbf{R}^2\). If \(x \in \mathbf{X}\), the tangent space \(\mathrm{T}_x(\mathbf{X})\) is identified with the line \(\mathbf{R}ix\) of \(\mathrm{T}_x(\mathbf{C}) = \mathbf{C}\), and the chosen orientation is the half-line \(\mathbf{R}_+ix\). If \(f\) is a continuous function on X, with values in a complex Banach space, the integral of \(f\) for the measure \(dz\) is given by the formula

\[
\int_ {\mathbb {X}} f (z) d z = i \rho \int_ {0} ^ {2 \pi} f (\rho e ^ {i \alpha}) e ^ {i \alpha} d \alpha .
\]

For example, if \(n\) is an integer, we have

\[
\int_ {\mathrm{X}} z ^ {n} d z = i \rho^ {n + 1} \int_ {0} ^ {2 \pi} e ^ {(n + 1) i \alpha} d \alpha = \left\{ \begin{array}{l l} 0 & \quad \text {if } n \neq - 1 \\ 2 i \pi & \quad \text {if } n = - 1. \end{array} \right.
\]

10.4.6 (“Complex case”). Let \(X^{c}\) be a separated complex analytic manifold, pure of dimension \(m\), and let \(\omega\) be a holomorphic differential form of degree \(m\) on \(X^{c}\). Let X be the real analytic manifold underlying \(X^{c}\); the form \(\omega\) identifies with a form of type \((m,0)\) on X (8.8.9); let \(\overline{\omega}\) be its conjugate (8.8.2). The form \(i^{m^{2}}\omega \wedge \overline{\omega}\) is a real differential form of degree \(2m\) on X; taking into account the canonical orientation of X (10.2.7), this form identifies with a measure on X, which is none other than the positive measure \(\operatorname{mod}(\omega)_{\mu}\) defined in 10.1.6, where \(\mu\) is the measure defined by the differential form \(i \, dz \wedge \overline{dz}\), in other words twice the Lebesgue measure \(\lambda^{\otimes 2}\) on \(\mathbf{C}\) (identified with \(\mathbf{R}^{2}\) in the usual way).

Let H be the space of holomorphic forms \(\omega\) of degree \(m\) on \(\mathbf{X}^c\) such that the measure \(\omega \wedge \overline{\omega}\) is bounded. For \(\alpha, \beta\) in H, the complex measure \(\alpha \wedge \overline{\beta}\) is bounded; setting

\[
(\alpha | \beta) = i ^ {m ^ {2}} \int_ {\mathbf {X}} \alpha \wedge \bar {\beta},
\]

one obtains a Hermitian form on H, which makes H into a Hilbert space.

